\section*{Capítulo \ref{ch:g4:world-physical-map} --- O mapa físico do mundo}

\begin{solution}{exo:g4:world-physical-map:1}
A \omterm{def:g4:world-physical-map:altitude}{altitude} de um lugar é a sua altura acima do \omterm{def:g4:world-physical-map:altitude}{nível do mar}, em metros.
Ela é medida a partir do \omterm{def:g4:world-physical-map:altitude}{nível do mar}, o nível médio da superfície do \omterm{def:g2:water-on-earth:ocean}{mar}.
\end{solution}

\begin{solution}{exo:g4:world-physical-map:2}
(a) Verde-escuro: entre 0 e 200 metros. (b) Amarelo-claro: entre 500 e
1\,000 metros. (c) Quase branco: acima de 4\,000 metros.
\end{solution}

\begin{solution}{exo:g4:world-physical-map:3}
(a) América do Sul. (b) Ásia. (c) África. (d) América do Norte.
\end{solution}

\begin{solution}{exo:g4:world-physical-map:4}
A \omterm{def:g1:landscapes-around-us:landform}{montanha} mais alta é o monte Everest (8\,849 metros). O maior \omterm{def:g1:landscapes-around-us:water}{lago} é o
\omterm{def:g2:water-on-earth:ocean}{mar} Cáspio; o \omterm{def:g1:landscapes-around-us:water}{lago} mais profundo é o \omterm{def:g1:landscapes-around-us:water}{lago} Baikal.
\end{solution}

\begin{solution}{exo:g4:world-physical-map:5}
Quanto mais escuro o azul, mais fundo o \omterm{def:g2:water-on-earth:ocean}{mar}: o \omterm{def:g2:water-on-earth:ocean}{mar} é raso perto das costas
(menos de 200 metros) e muito mais fundo no meio dos \omterm{def:g2:water-on-earth:ocean}{oceanos} (mais de
4\,000 metros em muitos lugares).
\end{solution}

\begin{solution}{exo:g4:world-physical-map:6}
O \omterm{def:g4:world-physical-map:physical-map}{mapa físico} mostra a \omterm{def:g4:world-physical-map:altitude}{altitude} da superfície. Na Antártida e na
Groenlândia, a superfície é o topo de uma camada grossa de gelo, que fica
bem acima do \omterm{def:g4:world-physical-map:altitude}{nível do mar}; por isso elas recebem as cores das terras altas.
\end{solution}

\begin{solution}{exo:g4:world-physical-map:7}
O fundo da Depressão Challenger (quase 11\,000 metros abaixo do \omterm{def:g4:world-physical-map:altitude}{nível do
mar}); o fundo médio do \omterm{def:g2:water-on-earth:ocean}{oceano} (cerca de 3\,700 metros abaixo); uma praia
(cerca de 0 metro); uma vila a 1\,500 metros; o topo do monte Everest
(8\,849 metros).
\end{solution}

\begin{solution}{exo:g4:world-physical-map:8}
Um \omterm{def:g1:landscapes-around-us:water}{rio} é formado por muitos riachos que se juntam, e as pessoas não
concordam sobre qual riacho é o seu verdadeiro começo. Medido a partir de
uma nascente mais distante, o \omterm{def:g1:landscapes-around-us:water}{rio} fica mais comprido. O Nilo e o Amazonas
têm comprimentos tão parecidos (cada um com cerca de 6\,700 a 6\,800
quilômetros) que a resposta depende dessa escolha.
\end{solution}

\begin{solution}{exo:g4:world-physical-map:9}
Um quarto de 30 milhões é sete milhões e meio de quilômetros quadrados (30
dividido por 4); um terço é 10 milhões (30 dividido por 3). O Saara, com
oito milhões e meio, é mais que um quarto da África, mas menos que um
terço.
\end{solution}

\begin{solution}{exo:g4:world-physical-map:10}
$2500 - 1000 = 1500$: o abrigo fica 1\,500 metros mais alto. O abrigo, entre
2\,000 e 4\,000 metros, é marrom; o \omterm{def:g1:landscapes-around-us:water}{lago}, a 1\,000 metros, fica no limite
entre o amarelo-claro (500 a 1\,000) e o marrom-claro (1\,000 a 2\,000):
qualquer um dos dois é aceito.
\end{solution}

\begin{solution}{exo:g4:world-physical-map:11}
10 quilômetros são 10\,000 metros. $10000 - 8849 = 1151$: o avião está mais
alto que o monte Everest, 1\,151 metros mais alto. Sobre a Depressão
Challenger, ele está a 10\,000 metros acima do \omterm{def:g2:water-on-earth:ocean}{mar}, e o fundo da fossa fica
a 11\,000 metros abaixo do \omterm{def:g2:water-on-earth:ocean}{mar}: $10000 + 11000 = 21000$ metros acima do
fundo da fossa.
\end{solution}
